\section{Proposed Scheme: A Counter-Based Search Index}
\label{sec:counter-index}

The paper's search has two problems. \textbf{The search key is permanent}: a server can keep $\dsk$ from one search
and test every later upload with it, so the scheme has no forward privacy. \textbf{Every search tries every item}:
the server tries $\dsk$ on all stored items and runs one lattice PRF chain per match. We replace the per-item tags
with an encrypted lookup table in the style of Cash et al.~\cite{cash2014dynamic}. The full design is in
\texttt{docs/counter-index-spec.md}.

\subsection{Design}

\paragraph{Keys.} The user derives three keys from $\Ku$, and none leaves the device: the keyword key
$K_w = \mathrm{HKDF}(\Ku, w)$, which is the paper's $\dsk$, and the pointer key $k_{ptr}$ and counter key $k_{ctr}$,
the two halves of one HKDF output.

\paragraph{Labels and rows.} The $j$-th file under keyword $w$ gets the label
$\mathit{label}(w, j) = \mathrm{PRF}(K_w, j)$, truncated to \CiLabelBytes. Its row stores the label with the entry
\[
\mathit{entry} = \mathrm{Enc}\big(k_{ptr},\ (\mathit{label},\ i,\ \mathrm{SHA3\text{-}256}(C))\big).
\]
The entry hides which file the row points to. The label inside stops a server from moving the entry to another
label, and the hash binds the row to one file. Each server keeps a dictionary per user from label to entry.

\paragraph{Counter table.} The table maps each keyword to $c_w$, the number of files under it. The user holds it
during a session. The servers keep a copy sealed under $k_{ctr}$, padded to 1024-byte blocks to hide the number of
keywords, versioned, and signed with $sk_u$, so a new device gets the same counters. A server accepts an upload only
with a valid signature and a higher version.

\subsection{Phases}

Register, Login steps L1 to L3 and the key update do not change, and since $\Ku$ survives a key update, so do the
labels.

\begin{protocol}{Login step L4 and Outsource with the counter index}
\begin{enumerate}[1.]
  \item \textbf{L4.} After L3, $U$ fetches the sealed table from each contacted server. It keeps the uploads whose
        signature verifies and whose table opens with the signed version inside, and takes the highest version.
  \item \textbf{O1.} For a file $d$ under $w$ with item id $i$, $U$ sets $j = c_w + 1$ and computes
        $\mathit{label}(w, j)$, $C = \mathrm{Enc}(\Ku, d)$ and the entry. It sends them with a signature to one
        server.
  \item \textbf{O2.} The server verifies the signature and stores the row and $C$, unless the label already exists
        or the user's $M$ slots are full.
  \item \textbf{O3.} After the server accepts, $U$ sets $c_w = j$ and uploads the new sealed table to all $N$
        servers. If the label already existed, an earlier upload was interrupted; $U$ moves to the next label.
\end{enumerate}
\end{protocol}

\begin{protocol}{Retrieve with the counter index}
\begin{enumerate}[R1.]
  \item If $c_w = 0$, $U$ returns nothing. Otherwise it sends the signed list $\mathit{label}(w, 1), \dots,
        \mathit{label}(w, c_w)$ to all $N$ servers.
  \item Each server verifies the signature, does $c_w$ dictionary lookups, and returns the rows it holds.
  \item $U$ drops rows whose label it did not ask for, opens each entry with $k_{ptr}$, drops entries whose inside
        label differs, and keeps duplicates once.
  \item $U$ sends a signed list of item ids to each server that returned hits and gets back the files.
  \item $U$ checks each file's hash against its entry and decrypts $d = \mathrm{Dec}(\Ku, C)$.
\end{enumerate}
\end{protocol}

Retrieve takes two round trips, because a server cannot follow an encrypted pointer. Unlike forward-private schemes
such as Bost's~\cite{bost2016sophos}, our keys come from a password through the TOPRF, so the user stores no key on
any device.
